% Source: Wolf (2012), Chapter 8, Equation (8.110), lines 1274--1293

\begin{theorem}[Detailed balance in the matrix representation]
    \label{thm:detailed_balance_transfer_matrix}
    \lean{transferMatrix_detailedBalance}
    \leanok
    This is the matrix-representation step in the proof of
    \cite[Chapter~8, Equation~(8.110)]{Wolf2012Quantum}.
    Let $T:M_d(\mathbb C)\to M_d(\mathbb C)$ be Hermiticity preserving and
    let $\Sigma:M_d(\mathbb C)\to M_d(\mathbb C)$ be linear.  If
    $\Sigma T^*=T\Sigma$, then
    \begin{align}
        \widehat\Sigma\,\widehat T^\dagger
        =\widehat T\,\widehat\Sigma.
    \end{align}
\end{theorem}

\begin{proof}
    \leanok
    Composition becomes matrix multiplication under the transfer-matrix
    representation.  Hermiticity preservation identifies the transfer matrix
    of $T^*$ with $\widehat T^\dagger$, so applying the representation to both
    sides of $\Sigma T^*=T\Sigma$ gives the displayed identity.
\end{proof}

\begin{theorem}[Matrix representation of the square-root sandwich]
    \label{thm:sigma_sandwich_transfer_matrix}
    \lean{transferMatrix_sigmaSandwich_eq_kronecker_and_posDef}
    \leanok
    This is the Kronecker identity used in the specialization of
    \cite[Chapter~8, Proposition ``Jordan condition number and detailed balance'']{Wolf2012Quantum}.
    If $\sigma>0$ and
    $\Sigma(X)=\sqrt\sigma\,X\sqrt\sigma$, then
    \begin{align}
        \widehat\Sigma
        =\overline{\sqrt\sigma}\otimes\sqrt\sigma>0.
    \end{align}
\end{theorem}

\begin{proof}
    \leanok
    Vectorizing conjugation by $\sqrt\sigma$ gives the displayed Kronecker
    product.  The square root is positive definite, and its entrywise
    conjugate is its transpose because it is Hermitian.  Transposition and
    Kronecker products preserve positive definiteness.
\end{proof}

\begin{theorem}[Hermitian representative from detailed balance]
    \label{thm:detailed_balance_hermitian_representative}
    \lean{Matrix.PosDef.inv_sqrt_mul_mul_sqrt_isHermitian_of_detailedBalance}
    \leanok
    This is the square-root conjugation in the proof of
    \cite[Chapter~8, Proposition ``Jordan condition number and detailed balance'']{Wolf2012Quantum}.
    Let $S,A\in M_n(\mathbb C)$, with $S>0$.  If
    $SA^\dagger=AS$, then
    $S^{-1/2}AS^{1/2}$ is Hermitian.
\end{theorem}

\begin{proof}
    \leanok
    Put $R=S^{1/2}$.  Positive definiteness makes $R$ invertible and
    Hermitian.  Multiplying $R^2A^\dagger=AR^2$ on the left and right by
    $R^{-1}$ gives
    $RA^\dagger R^{-1}=R^{-1}AR$, which says exactly that
    $R^{-1}AR$ is Hermitian.
\end{proof}

\begin{theorem}[The source diagonalizing similarity]
    \label{thm:detailed_balance_diagonalizing_similarity}
    \lean{Matrix.PosDef.exists_sqrt_mul_unitary_diagonalization_of_detailedBalance}
    \leanok
    This is the chosen change of basis in the proof of
    \cite[Chapter~8, Proposition ``Jordan condition number and detailed balance'']{Wolf2012Quantum}.
    Let $S,A\in M_n(\mathbb C)$, with $S>0$ and
    $SA^\dagger=AS$.  There are a unitary $U$ and a real diagonal matrix
    $D$ such that, for
    \begin{align}
        P=S^{1/2}U,
        \qquad
        Q=U^\dagger S^{-1/2},
    \end{align}
    one has
    \begin{align}
        PQ=QP=\Id,
        \qquad
        A=PDQ,
        \qquad
        \lVert P\rVert_\infty\lVert Q\rVert_\infty
        =\sqrt{\lVert S\rVert_\infty\lVert S^{-1}\rVert_\infty}.
    \end{align}
    This is a statement about Wolf's displayed, chosen diagonalizing basis;
    it does not identify that factor with the infimum $\kappa_T$.
\end{theorem}

\begin{proof}
    \leanok
    \uses{thm:detailed_balance_hermitian_representative}
    The preceding theorem makes $B=S^{-1/2}AS^{1/2}$ Hermitian.  Diagonalize
    it unitarily as $B=UDU^\dagger$ and conjugate back to obtain $A=PDQ$.
    Unitarity gives the two inverse identities and leaves the operator norm
    unchanged.  Finally the isometric functional calculus gives
    $\lVert S^{1/2}\rVert_\infty=\sqrt{\lVert S\rVert_\infty}$ and the
    corresponding identity for $S^{-1/2}$.
\end{proof}

\begin{theorem}[Fixed point from the detailed-balance sandwich]
    \label{thm:detailed_balance_sigma_fixed_point}
    \lean{fixedPoint_of_detailedBalance_sigmaSandwich}
    \leanok
    This is the fixed-point observation in
    \cite[Chapter~8, Proposition ``Jordan condition number and detailed balance'']{Wolf2012Quantum}.
    Let $\sigma>0$ and set
    $\Sigma(X)=\sqrt\sigma\,X\sqrt\sigma$.  If
    $\Sigma T^*=T\Sigma$ and $T^*(\Id)=\Id$, then $T(\sigma)=\sigma$.
\end{theorem}

\begin{proof}
    \leanok
    Evaluate detailed balance at the identity:
    $T(\sigma)=T(\Sigma(\Id))=\Sigma(T^*(\Id))=\Sigma(\Id)=\sigma$.
\end{proof}
